% GENERATED FILE. Edit canonical lessons, then run npm run generate:books.
% canonical-source-hash: fnv1a64:e2855e2c2546290b
% canonical-lessons: UR-C213-pahunchana, UR-C213-uthana, UR-C213-girana, UR-C213-phenkna, UR-C213-khinchna

\chapter{To deliver, To lift, To drop, To throw, To pull}
\label{ch:ur-to-deliver-to-lift-to-drop-to-throw-to-pull}
\hlchaptermodality{213}

\begin{chapteropening}
\textbf{By the end of this chapter:} \emph{I can say to deliver, to lift, to drop, to throw and to pull.}

Complete the last lesson of chapter 213: I can say to deliver, to lift, to drop, to throw and to pull.
\end{chapteropening}

\section[\texorpdfstring{\ur{پہنچانا} (pahuṅchānā)}{پہنچانا (pahuṅchānā)}]{\ur{پہنچانا} (pahuṅchānā) — to deliver, to take (somewhere)}

\label{lesson:UR-C213-pahunchana}



\begin{quote}
Before the new one: say the Urdu for to boil, then the Urdu for to taste.
\end{quote}

\subsection*{You'll want to know: \ur{پہنچانا}}
\textbf{\ur{پہنچانا}} — \emph{pahuṅchānā} — \textquotedblleft{}to deliver, to take (somewhere)\textquotedblright{}.

\begin{cousinweb}
\textbf{\ur{پہنچانا}} (\emph{pahuṅchānā}) is the causing form of \emph{pahuṅchnā}, \textquotedblleft{}to arrive\textquotedblright{}: to deliver is to make something arrive. CDIAL lists \emph{pahuṅchnā} under Sanskrit \emph{prabhūta}, \textquotedblleft{}much, great\textquotedblright{}, which Prakrit made into a verb, \textquotedblleft{}is strong, reaches\textquotedblright{} (CDIAL 8716).
\end{cousinweb}

\begin{grammarlens}[title={What you've built}]
One more word for this chapter.
\end{grammarlens}

\subsection*{Your turn}
\begin{itemize}
\raggedright
  \item \emph{Say it:} \emph{pahuṅchānā}
  \item \emph{Say it:} \emph{pahuṅchānā}, once more
  \item \emph{Say it:} say \emph{pahuṅchānā}
  \item \emph{Recall:} say the Urdu for to boil, then the Urdu for to taste, then say \emph{pahuṅchānā} again
\end{itemize}

\begin{tcolorbox}[breakable,colback=teal!4,colframe=teal!35!black,title={Before you move on}]
What does \ur{پہنچانا} mean? (\textquotedblleft{}To deliver, to take (somewhere)\textquotedblright{}.) Say it once more.
\end{tcolorbox}
\section[\texorpdfstring{\ur{اٹھانا} (uṭhānā)}{اٹھانا (uṭhānā)}]{\ur{اٹھانا} (uṭhānā) — to lift, to pick up}

\label{lesson:UR-C213-uthana}



\begin{quote}
Before the new one: say the Urdu for to taste, then the Urdu for to deliver.
\end{quote}

\subsection*{You'll want to know: \ur{اٹھانا}}
\textbf{\ur{اٹھانا}} — \emph{uṭhānā} — \textquotedblleft{}to lift, to pick up\textquotedblright{}.

\begin{cousinweb}
\textbf{\ur{اٹھانا}} (\emph{uṭhānā}) is a tadbhava, a word inherited and reshaped by sound change. It comes from \emph{*ut-sthāpayati}, \textquotedblleft{}raises\textquotedblright{}, a form CDIAL reconstructs from its descendants (CDIAL 1903). It is built on \emph{sthā-}, \textquotedblleft{}to stand\textquotedblright{}: to lift is to make something stand up.
\end{cousinweb}

\begin{grammarlens}[title={What you've built}]
One more word for this chapter.
\end{grammarlens}

\subsection*{Your turn}
\begin{itemize}
\raggedright
  \item \emph{Say it:} \emph{uṭhānā}
  \item \emph{Say it:} \emph{uṭhānā}, once more
  \item \emph{Say it:} say \emph{uṭhānā}
  \item \emph{Recall:} say the Urdu for to taste, then the Urdu for to deliver, then say \emph{uṭhānā} again
\end{itemize}

\begin{tcolorbox}[breakable,colback=teal!4,colframe=teal!35!black,title={Before you move on}]
What does \ur{اٹھانا} mean? (\textquotedblleft{}To lift, to pick up\textquotedblright{}.) Say it once more.
\end{tcolorbox}
\section[\texorpdfstring{\ur{گرانا} (girānā)}{گرانا (girānā)}]{\ur{گرانا} (girānā) — to drop}

\label{lesson:UR-C213-girana}



\begin{quote}
Before the new one: say the Urdu for to deliver, then the Urdu for to lift.
\end{quote}

\subsection*{You'll want to know: \ur{گرانا}}
\textbf{\ur{گرانا}} — \emph{girānā} — \textquotedblleft{}to drop\textquotedblright{}.

\begin{cousinweb}
\textbf{\ur{گرانا}} (\emph{girānā}) is the causing form of \emph{girnā}, \textquotedblleft{}to fall\textquotedblright{}: to drop is to make something fall. CDIAL lists \emph{girnā} under a reconstructed \emph{*girati}, \textquotedblleft{}drips, falls\textquotedblright{} (CDIAL 4159).
\end{cousinweb}

\begin{grammarlens}[title={What you've built}]
One more word for this chapter.
\end{grammarlens}

\subsection*{Your turn}
\begin{itemize}
\raggedright
  \item \emph{Say it:} \emph{girānā}
  \item \emph{Say it:} \emph{girānā}, once more
  \item \emph{Say it:} say \emph{girānā}
  \item \emph{Recall:} say the Urdu for to deliver, then the Urdu for to lift, then say \emph{girānā} again
\end{itemize}

\begin{tcolorbox}[breakable,colback=teal!4,colframe=teal!35!black,title={Before you move on}]
What does \ur{گرانا} mean? (\textquotedblleft{}To drop\textquotedblright{}.) Say it once more.
\end{tcolorbox}
\section[\texorpdfstring{\ur{پھینکنا} (pheṅknā)}{پھینکنا (pheṅknā)}]{\ur{پھینکنا} (pheṅknā) — to throw}

\label{lesson:UR-C213-phenkna}



\begin{quote}
Before the new one: say the Urdu for to lift, then the Urdu for to drop.
\end{quote}

\subsection*{You'll want to know: \ur{پھینکنا}}
\textbf{\ur{پھینکنا}} — \emph{pheṅknā} — \textquotedblleft{}to throw\textquotedblright{}.

\begin{cousinweb}
\textbf{\ur{پھینکنا}} (\emph{pheṅknā}) is a tadbhava, a word inherited and reshaped by sound change. It comes from \emph{*pheṅk}, \textquotedblleft{}throw\textquotedblright{}, a form CDIAL reconstructs from its descendants (CDIAL 9106). It is a form CDIAL gives beside a reconstructed \emph{*phekk}, \textquotedblleft{}throw\textquotedblright{}.
\end{cousinweb}

\begin{grammarlens}[title={What you've built}]
One more word for this chapter.
\end{grammarlens}

\subsection*{Your turn}
\begin{itemize}
\raggedright
  \item \emph{Say it:} \emph{pheṅknā}
  \item \emph{Say it:} \emph{pheṅknā}, once more
  \item \emph{Say it:} say \emph{pheṅknā}
  \item \emph{Recall:} say the Urdu for to lift, then the Urdu for to drop, then say \emph{pheṅknā} again
\end{itemize}

\begin{tcolorbox}[breakable,colback=teal!4,colframe=teal!35!black,title={Before you move on}]
What does \ur{پھینکنا} mean? (\textquotedblleft{}To throw\textquotedblright{}.) Say it once more.
\end{tcolorbox}
\section[\texorpdfstring{\ur{کھینچنا} (khīṅchnā)}{کھینچنا (khīṅchnā)}]{\ur{کھینچنا} (khīṅchnā) — to pull}

\label{lesson:UR-C213-khinchna}



\begin{quote}
Before the new one: say the Urdu for to drop, then the Urdu for to throw.
\end{quote}

\subsection*{You'll want to know: \ur{کھینچنا}}
\textbf{\ur{کھینچنا}} — \emph{khīṅchnā} — \textquotedblleft{}to pull\textquotedblright{}.

\begin{cousinweb}
\textbf{\ur{کھینچنا}} (\emph{khīṅchnā}) is a tadbhava, a word inherited and reshaped by sound change. It comes from \emph{*khiñc}, \textquotedblleft{}drag, pull\textquotedblright{}, a form CDIAL reconstructs from its descendants (CDIAL 3881).
\end{cousinweb}

\begin{grammarlens}[title={What you've built}]
One more word for this chapter.
\end{grammarlens}

\subsection*{Your turn}
\begin{itemize}
\raggedright
  \item \emph{Say it:} \emph{khīṅchnā}
  \item \emph{Say it:} \emph{khīṅchnā}, once more
  \item \emph{Say it:} say \emph{khīṅchnā}
  \item \emph{Recall:} say the Urdu for to drop, then the Urdu for to throw, then say \emph{khīṅchnā} again
\end{itemize}

\begin{tcolorbox}[breakable,colback=teal!4,colframe=teal!35!black,title={Before you move on}]
What does \ur{کھینچنا} mean? (\textquotedblleft{}To pull\textquotedblright{}.) Say it once more.
\end{tcolorbox}
